\chapter{Primer elemento constante de la responsabilidad: el daño y el perjuicio}
\section{Los elementos constantes de la responsabilidad}
Los elementos constantes de la responsabilidad son dos: el daño y el nexo causal. Se llaman «constantes» porque están presentes tanto en el régimen subjetivo como en el objetivo: siempre, sin excepción, se requiere acreditar el daño para que haya responsabilidad, y siempre se requiere acreditar el nexo causal.

\section{Daño y perjuicio}
\subsection{Definiciones}
\begin{definicion}\textbf{Daño:} es el hecho lesivo; la lesión a un bien jurídicamente amparado o protegido por el ordenamiento.\end{definicion}
\begin{definicion}\textbf{Perjuicio:} son las consecuencias materiales o inmateriales que el daño ocasiona al demandante. Estas consecuencias pueden ser de carácter patrimonial (material) o de carácter inmaterial.\end{definicion}
\begin{ejemplo}En un accidente de tránsito, el daño es el choque al vehículo (el hecho lesivo); el perjuicio es lo que cuesta volver a dejar las cosas en su estado inicial (la consecuencia del hecho lesivo).\end{ejemplo}
\subsection{Utilidad práctica de la distinción}
La distinción entre daño y perjuicio no es solo doctrinaria: tiene implicaciones prácticas. Las principales son dos:

\begin{itemize}
\item El estudio de la caducidad del medio de control.
\item La legitimación para demandar: la distinción entre víctimas directas e indirectas del daño.
\end{itemize}
\section{Primera utilidad: la caducidad}
\subsection{La regla general}
La caducidad se cuenta a partir del día siguiente de la ocurrencia del hecho (del daño), no desde que se consolidan los perjuicios. Hay situaciones en las que el daño y los perjuicios se consolidan en momentos distintos, pero esto de ninguna forma permite extender en el tiempo ni prolongar el término de caducidad: el término siempre se cuenta desde el día siguiente de la ocurrencia del daño. El hecho de que los perjuicios continúen en el tiempo no prorroga la caducidad.

\begin{comentario}La regla de caducidad de la reparación directa es: dos años contados a partir del día siguiente al de la ocurrencia de la acción u omisión causante del daño, o desde cuando el demandante tuvo o debió tener conocimiento del mismo si fue en fecha posterior y siempre que pruebe la imposibilidad de haberlo conocido en la fecha de su ocurrencia (ver Capítulo 8).\end{comentario}
\subsection{La regla del conocimiento del daño}
Hay casos particulares en los que la víctima no tiene conocimiento inmediato del daño o de su atribución al Estado. Por eso la norma prevé que, en situaciones excepcionales, el término se cuente desde que la víctima conoció o debió haber tenido conocimiento del daño.

\textbf{Filosofía de la regla.} La norma se pone en el papel del demandante: sin conocer el daño, o sin conocer que el daño se le debe atribuir al Estado, es imposible que presente la demanda. Esta regla es especialmente válida en materia médica, donde muchas veces la víctima no conoce inmediatamente el daño.

\begin{ejemplo}Transfusión de sangre: una persona es contagiada de VIH en una transfusión. El día de la transfusión no sabe que ha sido contagiada; puede tener ese conocimiento meses o un año después. Lo mismo ocurre con el cáncer: la persona no lo sabe desde el mismo día en que ocurre el hecho dañoso.\end{ejemplo}
\textbf{Carga de invocarla desde la demanda.} La regla del conocimiento debe plantearse con claridad desde la demanda. Lo que sucede hoy en día es que se presenta la demanda, se declara la caducidad y solo en el recurso el demandante invoca la regla del conocimiento. «Esa no es la forma.» La regla del conocimiento no es para salvarle el proceso a quien no la invocó: es para decir con claridad, desde la demanda, «conocí el daño en esta fecha, por esta razón».

\subsection{Casos prácticos de caducidad}
\paragraph{Primer caso: fechas de reingreso (desplazamiento)}
El primer caso involucra fechas de reingreso. Existe la tesis según la cual, en el desplazamiento, el daño es continuado: persiste en el tiempo y el término se cuenta cuando cesa, es decir, cuando la persona regresa o se restablece en condiciones similares a las que tenía, lo que permitiría demandar después del regreso. Es la posición mayoritaria, aunque es discutible que el daño se prorrogue en el tiempo (ver más adelante el debate sobre el desplazamiento forzado).

\paragraph{Segundo caso: el «falso positivo»}
En el segundo caso, los familiares recibieron una llamada en la que se les informaba de la muerte de la víctima, presentada oficialmente como baja en combate. Podría pensarse en aplicar la regla de la desaparición forzada (el término corre desde que aparece la víctima o desde la sentencia penal). Esa regla no es aplicable, por las siguientes razones:

\begin{itemize}
\item La regla especial es para la desaparición forzada: en ella los demandantes están en la incertidumbre, y por eso se espera a que aparezca la persona o a que haya una sentencia. Aplicarla a un caso de homicidio es extender una regla excepcional a otro supuesto de hecho.
\item En este caso, el daño es el homicidio y no la desaparición forzada; a los familiares se les comunicó la muerte.
\item Los demandantes denunciaron penalmente a los soldados. Si tuvieron conocimiento suficiente para acudir al proceso penal, también lo tenían para presentar la demanda de reparación directa en ese mismo momento. «Es extraño decir: para el penal sí tenía conocimiento en esa época, pero para la reparación directa, diez años después.» Hay una contradicción en esa argumentación.
\item No es cierto que haya que esperar al incidente de reparación en el proceso penal: quien inicia el proceso penal y espera, corre el riesgo de que opere la prescripción de la acción penal; por eso se inician simultáneamente ambos.
\end{itemize}
Hay quienes defienden otra lectura: que la sentencia sobre flexibilización probatoria en estos casos permitía contar el término desde la sentencia, que la Corte Constitucional ha dicho que la desaparición forzada solo cesa cuando aparece la persona, que en el caso no aparecieron los restos y que, si los familiares acudieran a la Unidad de Búsqueda de Personas Desaparecidas, allí lo seguirían buscando (por lo que la desaparición seguiría siendo permanente), y que una denuncia penal (cuyo deber de investigar recae en la Fiscalía) no equivale a tener los mínimos para demostrar el daño y el nexo causal en una demanda de reparación directa. Ninguna de las posturas es descabellada, y es probable que, por las reglas de flexibilización en materia de derechos humanos, un juez concluya que la plena certeza se tuvo con la sentencia; pero esa conclusión es muy controvertible: no tiene sentido que se pueda activar al juez penal y no al juez contencioso.

\begin{comentario}Crítica: hay una tendencia a flexibilizar reglas en casos de graves violaciones de derechos humanos en la que «se les olvida el derecho procesal». Las defensas procesales en materia de derechos humanos «están de moda», pero no pueden llevar a olvidarse del derecho procesal.\end{comentario}
\paragraph{Tercer caso: la captura}
En el tercer caso se discutió desde cuándo se cuenta la caducidad cuando hay una captura: si desde la declaratoria de ilegalidad de la captura o desde que la persona queda en libertad. El problema es de antijuridicidad del daño: ¿en qué momento el daño es antijurídico?

\begin{itemize}
\item Una postura sostuvo que el daño es antijurídico desde que se declara la ilegalidad de la captura, porque se vulneraron los derechos de la persona, independientemente de que después se le declare penalmente responsable.
\item Sin embargo, la detención en sí misma no es necesariamente antijurídica, por más ilegal que haya sido la captura: si la persona termina condenada (por ejemplo, a muchos años de cárcel), el tiempo que estuvo detenida se le descuenta de la pena. Solo cuando se la absuelve o se declara la prescripción de la acción penal se sabe realmente que tenemos un daño antijurídico, porque la persona nunca debió haber estado privada de la libertad.
\item Sobre el momento exacto, se cuenta desde lo último que ocurra entre la ejecutoria de la absolución y la libertad: si a la persona la absuelven y le dan la libertad cuatro o cinco días después, sería un error contar desde la absolución, porque esos días quedarían sin poder reclamarse; el término debe correr desde que la persona tiene derecho real a reclamar la totalidad del daño.
\end{itemize}
\paragraph{Cuarto caso: la obra de Transmilenio y el predio}
En el cuarto caso, la obra de Transmilenio afectó un predio. El daño se consolida cuando termina la obra: solo en ese momento se puede determinar qué tipo de daño hubo. Por ejemplo, el contratista puede usar el predio para almacenar ladrillos y material de obra durante un tiempo y luego retirarlo (ocupación temporal), o puede ocurrir que, al terminar la obra, el propietario pierda definitivamente parte de su propiedad (ocupación permanente). A partir de ese momento se cuenta la caducidad, aunque los perjuicios sigan causándose después.

\begin{ejemplo}¿Qué ordena el juez si comprueba la ocupación? En la reparación directa se ordena reconocer al propietario el valor del predio. La sentencia condenatoria se vuelve un título: le pagan, por ejemplo, 100 millones de pesos, queda reparado integralmente y con esa sentencia se traslada la propiedad a favor de la entidad (en últimas, el Estado termina comprando el predio). No se ordena destruir la obra, porque sería más gravoso para el interés público. Otra cosa es que, antes o durante el proceso, se soliciten medidas cautelares para suspender la obra.\end{ejemplo}
\subsection{Caducidad, lesa humanidad y desplazamiento forzado}
Durante muchos años el Consejo de Estado sostuvo que, como la acción penal por delitos de lesa humanidad (por ejemplo, la desaparición forzada) no prescribe, tampoco caducaba la reparación directa. Esa tesis se abandonó en la sentencia de unificación de 2019 sobre caducidad y lesa humanidad (sentencia de unificación del 29 de enero de 2020, expediente 61033): el hecho de que la acción penal sea imprescriptible no quiere decir que no caduque la reparación directa. Lo único que interesa para la caducidad es la existencia de barreras que impidan el acceso a la administración de justicia: si el demandante tiene barreras que le impiden acudir al juez administrativo, eso permite extender el término.

\begin{ejemplo}Barreras: amenazas, estar secuestrado, estar en el exilio (asilo); cualquier situación que impida a la persona acudir al juez.\end{ejemplo}
\textbf{El debate actual sobre el desplazamiento forzado.} Estas discusiones volvieron a surgir porque había consejeros que contaban la caducidad desde que ocurre el desplazamiento y otros que no. Se presentaron varias tutelas y la Corte Constitucional realizó una audiencia cuyas preguntas parecen orientadas a volver a la teoría anterior.

\begin{comentario}Crítica: volver a esa teoría sería grave, porque podría revivir hoy demandas por desplazamientos de los años 70 en las bananeras; el daño no se prorroga en el tiempo. Además, saber cuándo la persona volvió a tener el mismo estatus social o el mismo ingreso que tenía antes del desplazamiento es «absurdamente» difícil de sostener, y la caducidad tiene que ser clara para que haya seguridad jurídica. La jurisprudencia está generando así dos problemas: (i) la inaplicación del derecho procesal y (ii) la confusión entre el concepto de daño y la violación de derechos («hubo un daño porque se violó el derecho a la dignidad, hubo otro porque se violó el derecho a la vida…»). No se puede decir que hay distintos daños por cada derecho; se está aplicando mucho lenguaje constitucional y de derechos fundamentales a una disciplina que no necesariamente es para eso. Una cosa es proteger a la víctima y otra es el régimen de responsabilidad.\end{comentario}
\textbf{Temporalidad: ocupación temporal y permanente.} El mismo problema se presenta con la ocupación: solo al finalizar la obra se sabe si el daño es temporal o permanente. En el desplazamiento ocurre algo parecido con los perjuicios: no es lo mismo desplazarse por dos semanas (se pierde un poco del cultivo) que por tres o seis meses.

\section{Segunda utilidad: víctimas directas e indirectas}
La distinción entre daño y perjuicio permite diferenciar, dentro del proceso de reparación directa, entre:

\begin{definicion}\textbf{Víctima directa:} la persona que padece el daño y padece los perjuicios.\end{definicion}
\begin{definicion}\textbf{Víctimas indirectas:} quienes no padecen el daño, pero sí sufren perjuicios.\end{definicion}
\begin{ejemplo}En un caso de privación de la libertad, el privado de la libertad es la víctima directa; sus familiares son víctimas indirectas. La distinción tiene repercusiones prácticas en la cuantificación de los perjuicios.\end{ejemplo}
\section{Características (requisitos) del perjuicio indemnizable}
Para que un perjuicio pueda ser indemnizado en un proceso judicial debe cumplir cuatro requisitos:

\begin{itemize}
\item Ser \textbf{cierto}.
\item Ser \textbf{personal}.
\item Ser \textbf{lícito}.
\item Ser \textbf{subsistente} (debe subsistir al momento de la sentencia).
\end{itemize}
\subsection{El perjuicio cierto}
\begin{definicion}\textbf{Perjuicio cierto:} es aquel que no es meramente eventual o hipotético.\end{definicion}
\begin{ejemplo}Freddy, el futbolista. Freddy sale de la universidad y tiene un accidente de tránsito (el demandado puede ser un particular o el Estado). Tiene una incapacidad médica de 30 días que no le permite trabajar. Como es contratista, su lucro cesante cierto es el monto de los honorarios de su contrato de prestación de servicios, calculado por el término de la incapacidad. En cambio, Freddy no podría reclamar que, como todos los domingos es «el Messi del barrio», de no haber ocurrido el accidente estaría jugando en el Real Madrid ganando lo que gana Mbappé: ese perjuicio es meramente eventual, hipotético; existe en la mente de Freddy, pero no tiene ninguna base real de certeza.\end{ejemplo}
\textbf{a) Certeza no es lo mismo que actualidad.} Que el perjuicio sea cierto no significa que deba estar consolidado al presentar la demanda. Es procedente reclamar perjuicios futuros que son ciertos pero que aún no se han consolidado. El perjuicio puede ser pasado, presente o futuro.

\begin{ejemplo}Prótesis. Un soldado entra a un campo minado que no está debidamente señalizado, pisa una mina y pierde una pierna; debe usar una prótesis. Las prótesis tienen una vida útil: supongamos que deben cambiarse cada cinco años. Si el soldado tiene una expectativa de vida de 45 años, se calcula cuántas prótesis tendrá que comprar y ese gasto, que todavía no se ha hecho, se puede reclamar porque hay certeza de que se va a realizar. El reto es cómo calcularlo (por ejemplo, con el precio actual de la prótesis): lo ya pagado se indexa al momento de la sentencia; lo futuro no puede indexarse, pero sí se puede obtener su valor.\end{ejemplo}
\begin{ejemplo}Junta médico laboral. Un soldado sufre un accidente y está en proceso de calificación de invalidez, pero aún no es claro el porcentaje de pérdida de capacidad laboral. ¿Es un problema de daño o de perjuicio? La jurisprudencia ha concluido que la Junta Médico Laboral cuantifica el perjuicio, pero no el daño. En el caso del soldado que pisa la mina, el daño (la lesión a su integridad corporal) se consolidó cuando pisó la mina y perdió la pierna. Por eso debe demandar dentro de los dos años siguientes a la ocurrencia del daño; el dictamen de la Junta podrá aportarse como prueba sobreviniente si se expide durante el proceso (explicando que no podía aportarse con la demanda porque no existía). Ni siquiera es necesario ese dictamen para ganar: se puede pedir un dictamen pericial en el proceso para cuantificar el perjuicio.\end{ejemplo}
\begin{ejemplo}Menores de edad y cirugías futuras. Los menores lesionados a veces no pueden ser operados de inmediato y se sabe que serán sometidos a una cirugía dentro de cinco u ocho años. En estos casos se recomienda formular las pretensiones previendo que la cuantía se ajustará a lo que se demuestre en el proceso y la posibilidad de aportar dictámenes que se produzcan durante el trámite.\end{ejemplo}
\textbf{b) La certeza recae sobre el perjuicio, no sobre su cuantía.} Tanto en lo contencioso administrativo como en lo civil, los códigos prevén que cuando se demuestra el daño pero no la cuantía del perjuicio, el juez dicta una sentencia condenatoria en abstracto, que fija unas bases para que luego el demandante promueva un incidente de liquidación de perjuicios (un pequeño proceso dentro del mismo proceso), en el que solo se discute la cuantificación.

\begin{ejemplo}El vehículo embargado. En un proceso ejecutivo embargan y secuestran el vehículo del demandado. Cuando se paga la deuda y se lo devuelven, el vehículo está desvalijado porque no se ejerció su custodia. El dueño demanda, prueba el daño (tenía un buen vehículo y se lo devolvieron destruido), pero no prueba cuánto valía el motor, el radio o la pintura robados. El juez puede condenar porque está demostrado el daño, pero debe hacerlo en abstracto: la sentencia dirá, por ejemplo, que el daño emergente se liquidará en un incidente teniendo en cuenta que el vehículo era de tal marca, el motor de tal marca, la pintura de tal color, etc.\end{ejemplo}
\textbf{c) Certeza del nexo entre el daño y el perjuicio.} La certeza también se predica de la relación que debe existir entre el daño alegado en la demanda y el perjuicio reclamado en las pretensiones. Muchas demandas no demuestran esa relación, lo que lleva al juez a desestimar esos perjuicios.

\begin{ejemplo}Privación injusta de la libertad. A Pedro lo privan de la libertad y luego lo absuelven porque se demostró que el delito lo cometió otra persona. Son perjuicios ciertos: lo que dejó de trabajar, lo que pagó al abogado de la defensa y los gastos de transporte de la familia para visitarlo (si no hubiera estado privado de la libertad, la familia no habría tenido que hacer esos gastos). En cambio, el gasto en útiles de limpieza personal que usó en el centro de reclusión, aunque sea un gasto que efectivamente se causó, no tiene relación de certeza con la privación: no se deriva realmente del hecho dañoso (no es que la persona haya empezado a asearse por estar privada de la libertad).\end{ejemplo}
\subsection{Un punto intermedio: la pérdida de oportunidad}
Entre los perjuicios ciertos y los meramente eventuales o hipotéticos se ubica la pérdida de oportunidad. Es un problema complejo, porque ni en Colombia ni en el derecho comparado hay acuerdo sobre si se trata de un problema del perjuicio o del nexo causal: algunos la tratan como problema del perjuicio, otros como problema del nexo causal, y hay países que la usan para ambas cosas. En Colombia hay jurisprudencia en los dos sentidos. Los franceses la utilizan para el nexo causal.

\textbf{a) Enfoque desde el perjuicio.} El primer caso que utilizó la figura en el ámbito de los perjuicios fue un caso del \emph{common law}, de principios de 1900:

\begin{ejemplo}El concurso de belleza. Una empresa organiza un concurso de belleza; se inscriben 100 personas y el organizador escoge cinco finalistas, a quienes envía por correo físico la notificación para presentarse a la gran final en determinada fecha y ciudad. El premio es de 100.000 dólares. A cuatro finalistas la notificación les llega una semana antes; a la quinta le llega el día siguiente a la final, por lo que nunca pudo concursar. La quinta finalista demanda por pérdida de oportunidad y por primera vez se reconoce ese perjuicio. No se sabe si ella iba a ganar, pero sí que tenía una oportunidad cierta y seria de hacerlo: estaba entre las cinco finalistas, no en el puesto 99.\end{ejemplo}
\textbf{Requisitos de la pérdida de oportunidad:}

\begin{itemize}
\item Una oportunidad cierta de obtener un beneficio (probabilidad de obtener un provecho).
\item Una oportunidad seria.
\item Que esa oportunidad se frustre como consecuencia del hecho dañoso.
\item Que la frustración sea definitiva: el provecho no puede obtenerse de ninguna otra forma.
\end{itemize}
\textbf{Tasación.} Lo que sigue es la discusión de cómo se tasa: ¿hay que darle a la quinta finalista los 100.000 dólares o una parte? Es un tema de arbitrio judicial y de equidad. Una forma de hacerlo es dividir el premio entre cinco: la oportunidad perdida valía una quinta parte, 20.000 dólares. La consecuencia de este enfoque es una tasación porcentual del perjuicio.

\textbf{b) Enfoque desde el nexo causal (materia médica).}

\begin{ejemplo}A una persona le disparan tres veces (en la pierna, en el pulmón y en la mano). Llega al hospital y no la atienden de inmediato, sino dos horas después, y muere. ¿Hasta qué punto la demora contribuyó a la muerte o disminuyó la probabilidad de recuperación? Nunca se podrá determinar: quizás iba a morir de cualquier manera porque venía con lesiones muy graves. Aquí se ha usado la pérdida de oportunidad para preguntarse hasta qué punto la falla del servicio médico contribuyó a impedir que la persona sobreviviera.\end{ejemplo}
\begin{comentario}Valoración: el enfoque más adecuado es el del perjuicio. En materia causal no parece adecuado hablar de pérdida de oportunidad: realmente es un tema de prueba indiciaria y de estándares probatorios (por ejemplo, discutir si en materia médica el estándar de causalidad debe ser de probabilidad preponderante y no absoluta). Con todo, es un tema con discusiones muy amplias; el riesgo de la figura es la confusión entre daño y causalidad.\end{comentario}
\subsection{El perjuicio personal}
El carácter personal se relaciona con el tipo de obligación que nace de la responsabilidad extracontractual: obligaciones de carácter personal y no de carácter real. Esto significa que el único que puede reclamar un perjuicio es quien lo padeció.

\begin{itemize}
\item Quien reclama es quien sufre, no el daño, sino el perjuicio. Por ejemplo, el hijo puede reclamar perjuicios morales por la muerte de su padre porque es él quien padece el fallecimiento de su ser querido.
\item Por eso hay que distinguir con toda claridad en qué condición actúa cada demandante: puede reclamar perjuicios propios (\emph{iure proprio}) o perjuicios en calidad de heredero (\emph{iure hereditario}). Ambas reclamaciones son procesalmente válidas, pero deben explicarse desde la demanda: hay que alegar y acreditar la calidad de heredero, además de probar el perjuicio.
\end{itemize}
\begin{ejemplo}Pedro privado de la libertad. Pedro es privado de la libertad, es absuelto y tiene dos años para demandar.\end{ejemplo}
\begin{itemize}
\item \textbf{Escenario 1 – Pedro está vivo.} Demandan Pedro y sus familiares (A, B y C). Pedro reclama sus perjuicios personales (lucro cesante, perjuicios morales) y sus familiares reclaman sus propios perjuicios por la privación de su pariente. Todos actúan por derecho propio.
\item \textbf{Escenario 2 – Pedro muere antes de la demanda.} Demandan solo A, B y C. Reclaman sus perjuicios morales propios, pero además los hijos pueden actuar como herederos de Pedro y reclamar el lucro cesante que sufrió Pedro, el daño emergente que pagó Pedro e incluso los perjuicios morales de Pedro. Esto hay que especificarlo en la demanda; muchas veces se olvida y en la sentencia solo se reconocen los perjuicios propios.
\item \textbf{Escenario 3 – Pedro muere durante el proceso.} Es un tema netamente procesal: hay sucesión procesal. Lo que hacen los jueces es decretar la condena por lo reclamado por Pedro a favor de la masa sucesoral, o, si la sucesión ya terminó, a favor de quien corresponda, para que en el marco de la sucesión se adjudique ese crédito.
\end{itemize}
\textbf{Daños a las cosas.} El Código Civil establece que, cuando se trata de daños a las cosas, está legitimado en la causa por activa el propietario, poseedor o tenedor del bien al momento de la ocurrencia del daño. Debe demostrarse el derecho sobre la cosa, pero el derecho a reclamar no es un derecho real, sino personal.

\begin{ejemplo}En el proceso ejecutivo de A contra B se embarga y secuestra el vehículo de B, recién salido del concesionario. Cuando se lo devuelven, el vehículo está absolutamente destruido. Después B vende el carro a C. ¿Quién puede demandar al Estado? B, porque era el propietario cuando se causó el daño. C no puede demandar al Estado, porque la obligación de responsabilidad extracontractual es personal y no real: por el hecho de vender el carro no se transfiere el derecho a demandar. Lo que hay entre B y C es un problema contractual (C, a lo sumo, podría demandar a B por vicios ocultos).\end{ejemplo}
\subsection{El perjuicio lícito}
\begin{definicion}\textbf{Perjuicio lícito:} aquel que es consecuencia de una afectación a un bien o interés amparado o protegido por el ordenamiento jurídico.\end{definicion}
\textbf{No confundir licitud con titularidad.} La licitud del perjuicio no tiene nada que ver con la titularidad del derecho. En el ejemplo anterior, si C demanda, a veces los tribunales dicen que el perjuicio no es lícito; eso es un error: el perjuicio sí es lícito, el problema es que C no es el titular del derecho, que es un problema del carácter personal del perjuicio.

\begin{ejemplo}Hipótesis (hechos supuestos, con fines ilustrativos). Un criminal conocido es dado de baja estando ya esposado e inmovilizado, por un policía que actuó por rabia. Su hijo demanda al Estado. Puede reclamar los perjuicios morales por la muerte de su padre y el daño emergente por los gastos funerarios. Pero no puede reclamar, como lucro cesante, las utilidades que su padre dejó de recibir por el tráfico de cocaína: es un perjuicio absolutamente ilícito, porque la venta de cocaína no está permitida.\end{ejemplo}
\begin{ejemplo}El caso de alias «Bonais» (caso real). La policía persigue a un hombre conocido con ese alias, al parecer armado, y lo mata. La familia reclama y se condena al Estado (no lo podían ejecutar). Uno de los puntos de mayor discusión en la Sala fue si había que reconocer lucro cesante, porque nunca quedó claro si era microtraficante o vendedor informal: había pruebas de que vendía de manera informal y no se demostró lo contrario. Se aplicó entonces la presunción de actividad económica: como se demostró que tenía una actividad, pero no su valor, se liquidó con el salario mínimo (ver reglas de la renta, sección 4.8).\end{ejemplo}
\subsection{El perjuicio subsistente}
\begin{definicion}\textbf{Perjuicio subsistente:} el perjuicio debe seguir existiendo al momento del fallo; es decir, que no haya sido reparado con antelación y plenamente.\end{definicion}
La lógica es clara: si cuando se va a dictar sentencia la víctima ya obtuvo la indemnización del perjuicio, el juez no puede volver a darle una indemnización.

\begin{ejemplo}En un accidente de tránsito, la víctima demanda al conductor y durante el proceso la aseguradora le paga. El juez no puede condenar por la totalidad de lo reclamado, porque parte ya fue pagada: el perjuicio no es subsistente, o lo es solo parcialmente (lo que no cubrió la aseguradora).\end{ejemplo}
\textbf{La \emph{compensatio lucri cum damno}.} Un autor italiano, De Cupis, escribió un libro sobre el daño en el que planteó un problema relacionado con la subsistencia del perjuicio: el principio de la \emph{compensatio lucri cum damno}, definido como «la reducción del monto del daño resarcible por la concurrencia del lucro». Significa que a partir del daño el demandante no puede obtener lucro o beneficio: si la aseguradora ya pagó, reconocer la totalidad le generaría un enriquecimiento sin causa real.

El origen del problema es la coexistencia de la reparación judicial del daño con su compensación por otro mecanismo no judicial distinto al directamente responsable (seguridad social, solidaridad, reparación administrativa): la responsabilidad no busca lucrar a la víctima. Las preguntas que plantea son:

\begin{itemize}
\item ¿Es posible el cúmulo de indemnizaciones?
\item ¿Debe descontarse lo recibido en virtud del mecanismo no judicial de la reparación integral ordenada en sede judicial?
\item ¿El tercero que realiza el pago de la compensación no judicial puede repetir contra el responsable (acción subrogatoria)?
\end{itemize}
\textbf{Argumentos que se confrontan:}

{\small\begin{xltabular}{\linewidth}{|X|X|}\hline
\rowcolor{navy!12}\textbf{Argumentos para permitir la acumulación} & \textbf{Argumentos para descontar} \\\hline\endhead
Tienen una causa distinta (indemnización generada por responsabilidad patrimonial / seguridad social reglada por normas del derecho laboral) & Mismo hecho generador del daño \\\hline
La compensación no judicial no se calcula con base en el daño, ni lo disminuye & Evitar una doble indemnización \\\hline
Al tener fuente distinta, no se genera un enriquecimiento sin causa & Pueden coexistir porque la compensación no judicial es tarifada y la reparación judicial es integral, pero se deben aplicar los descuentos \\\hline
\end{xltabular}}
\textbf{Dos casos para analizar:}

\begin{ejemplo}Caso 1 – El soldado y la pensión de sobrevivientes. El soldado pisa una mina en un campo no señalizado y muere. Su familia demanda al Ejército y reclama daño emergente, lucro cesante (los salarios por su expectativa de vida) y perjuicios morales. Paralelamente, a la cónyuge (también demandante) le reconocen una pensión de sobrevivientes. ¿Debe descontarse la pensión del lucro cesante?\end{ejemplo}
\begin{ejemplo}Caso 2 – El maestro y la ARL. A un maestro de una escuela pública se le cae encima el tablero mientras dicta clase, porque la escuela no hizo el mantenimiento adecuado; sufre un accidente de trabajo. Durante la reparación directa, la ARL le reconoce una suma (por ejemplo, 10 salarios) por el accidente. ¿Debe descontarse de la indemnización?\end{ejemplo}
Una primera respuesta es que no debe descontarse en ninguno de los dos casos (son causas y fuentes de obligación distintas; en el caso de la ARL se asegura un riesgo de manera objetiva). En contra, puede sostenerse que en el primer caso lo que se cubre es exactamente lo mismo.

\begin{comentario}Argumento a favor del descuento: el hecho dañoso es el mismo; con la pensión se está reconociendo, por vía del derecho laboral, parte del sueldo que recibía el soldado, y en el proceso también se reclama ese sueldo; por más que tengan naturaleza jurídica distinta, resulta redundante.\end{comentario}
\textbf{La sentencia de unificación 66841 (pensión de invalidez y de sobrevivientes de miembros de la fuerza pública).} Esta sentencia de unificación resolvió que no hay que descontar. Sus criterios son:

{\small\begin{xltabular}{\linewidth}{|X|X|X|}\hline
\rowcolor{navy!12}\textbf{Criterio} & \textbf{Lucro cesante} & \textbf{Pensión} \\\hline\endhead
Causa & Falla del servicio & Riesgos laborales inherentes a la prestación del servicio \\\hline
Fuente jurídica & Art. 90 de la C.P. & Sistema de seguridad social (solidaridad) \\\hline
Naturaleza del pago & Indemnizatorio & Contraprestación social (laboral) no indemnizatoria \\\hline
Extinción de la obligación indemnizatoria & Sí & No \\\hline
\end{xltabular}}
\begin{itemize}
\item \textbf{Subrogación:} no existe subrogación legal a favor de quien otorga la pensión. Se distinguen estos casos de aquellos en los que el Estado efectúa cotizaciones al sistema de seguridad social para trasladar los riesgos profesionales a una entidad diferente que los asegura (ISS, CAJANAL, ARP o ARL).
\item \textbf{Conclusión:} la acumulación no genera enriquecimiento sin causa de la víctima.
\end{itemize}
\begin{comentario}Crítica a la sentencia: el criterio de la «extinción de la obligación» no es convincente, porque es precisamente lo que había que discutir (si pagar la pensión extingue o no la obligación); más que un criterio, es la conclusión. Además, la unificación solo se dio para pensiones de invalidez y de sobrevivientes de la fuerza pública.\end{comentario}
\textbf{El criterio decisivo: la subrogación.} Lo realmente importante de la sentencia es la subrogación: si hay subrogación, hay que descontar; si no la hay, no. En el caso de las ARL hay norma especial de subrogación: la ARL puede subrogarse y después demandar al responsable para recuperar lo pagado. Como hay subrogación, el pago sí extingue la obligación y recibir dos veces generaría un enriquecimiento a partir del hecho dañoso. En los casos de ARL todavía no hay unificación, pero con el criterio de esa sentencia ya se puede entender hacia dónde va la jurisprudencia en accidentes de trabajo.

\begin{comentario}Otros argumentos. Que el Ejército tenga sus propias cajas y administradores de pensiones sería otro argumento para descontar, pero la unificación dice lo contrario y debe respetarse. Por otra parte, según la línea de la Sala de Casación Laboral de la Corte Suprema de Justicia, en la indemnización por culpa patronal de un trabajador del sector privado no se descuenta lo que ya pagó la ARL por pérdida de capacidad laboral (por ejemplo, si se condena al empleador a 100 salarios y la ARL ya pagó 20 según la tabla, no se descuentan esos 20).\end{comentario}
\textbf{Lucro cesante por lesiones y la incapacidad.} ¿Puede Freddy (contratista con incapacidad de dos meses, pagada al 66 \% de su cotización por ser accidente común) reclamar la totalidad de sus honorarios? El Consejo de Estado, en lucro cesante por lesiones corporales, ha reconocido un porcentaje equivalente al de la incapacidad, con dos justificaciones:

\begin{itemize}
\item La lesión restringe la autonomía de la persona en cuanto al ejercicio de su profesión. Se citó una sentencia de los años 90 de un torero que recibe una cornada por una falla del servicio, deja de ser torero y tiene que emplearse en otra actividad.
\item La segunda razón, más potente: al estar incapacitada, la persona tiene que hacer un doble esfuerzo para hacer lo mismo que hacen los demás; por eso se le reconoce el porcentaje de la incapacidad.
\end{itemize}
\section{Tipologías de perjuicios en la jurisdicción de lo contencioso administrativo}
Las tipologías de perjuicios de la jurisdicción de lo contencioso administrativo son distintas a las de la jurisdicción ordinaria civil. Aquí se expone la posición del Consejo de Estado, por ser la que corresponde a la responsabilidad del Estado, con algunas diferencias frente a la ordinaria. Recientemente, una sentencia importante de la Corte Suprema de Justicia intentó acercar la jurisdicción ordinaria a la contenciosa y acogió varios criterios de esta, pero siguen siendo disímiles. Ambas coinciden en el daño emergente, el lucro cesante y los perjuicios morales; de ahí en adelante las tipologías son distintas.

\begin{definicion}\textbf{Perjuicios materiales (patrimoniales):} los que se reflejan directamente en el patrimonio del demandante.\end{definicion}
\begin{definicion}\textbf{Perjuicios inmateriales:} los que van más allá del patrimonio del demandante y afectan aspectos morales, espirituales, etc.\end{definicion}
{\small\begin{xltabular}{\linewidth}{|X|X|}\hline
\rowcolor{navy!12}\textbf{Categoría} & \textbf{Tipología} \\\hline\endhead
Patrimoniales o materiales & Daño emergente \\\hline
Patrimoniales o materiales & Lucro cesante \\\hline
Inmateriales & Perjuicios morales \\\hline
Inmateriales & Daño a la salud \\\hline
Inmateriales & Daños por la vulneración de bienes constitucional y convencionalmente protegidos \\\hline
\end{xltabular}}
\section{El daño emergente}
\begin{cita}«Art. 1614. Entiéndese por daño emergente el perjuicio o la pérdida que proviene de no haberse cumplido la obligación o de haberse cumplido imperfectamente, o de haberse retardado su cumplimiento (…)» (Código Civil)\end{cita}
\begin{definicion}\textbf{Daño emergente:} la disminución del patrimonio; la pérdida patrimonial que sufre el demandante a raíz del hecho dañoso. El patrimonio era uno antes del daño y, al ocurrir el daño, baja.\end{definicion}
El daño emergente puede ser:

\begin{itemize}
\item \textbf{Consolidado (pasado):} cuando se presenta la demanda ya se sufrió la pérdida patrimonial. Debe actualizarse a valor presente.
\item \textbf{Futuro:} cuando se presenta la demanda aún no se ha sufrido la pérdida (el caso de las prótesis o de la operación del niño).
\end{itemize}
\subsection{La actualización (indexación)}
La fórmula de actualización no es solo del daño emergente: se aplica siempre para actualizar los perjuicios materiales. Actualizar es indexar, es decir, traer a valor presente, porque el dinero pierde valor con el paso del tiempo. Donde es más evidente y sencillo de ver es en el daño emergente.

\formula{Va = Vh × (Índice final ÷ Índice inicial)}
\begin{itemize}
\item \textbf{Va:} valor actualizado.
\item \textbf{Vh (valor histórico):} el valor del daño al momento de su causación.
\item \textbf{Índice final:} IPC del mes de la presentación de la demanda (o del mes de la sentencia, si se liquida la condena).
\item \textbf{Índice inicial:} IPC del mes de la ocurrencia del daño (o de la fecha en que se hizo la erogación de dinero).
\end{itemize}
\textbf{¿Dónde se consulta el IPC?} En la página del DANE hay un cuadro de Excel con los índices históricos del IPC (Índice de Precios al Consumidor). Se toma el índice del mes para el cual se quiere hacer la actualización y se divide por el índice vigente a la fecha en que se hizo la erogación; el resultado se multiplica por el valor que se quiere actualizar.

\subsection{¿Para qué se calculan los perjuicios? La línea del tiempo del proceso}
La liquidación de perjuicios no es un problema solo del juez: es un problema del juez y de las partes. El abogado (del demandante o del demandado) tiene que verificar que la liquidación esté bien hecha, porque quizás deba apelarla («señor juez, usted me dio la mitad de lo que me tenían que dar, porque liquidó mal»). Para entender esto hay que tener en cuenta la línea del tiempo de un proceso de responsabilidad extracontractual del Estado:

\begin{itemize}
\item Ocurre el hecho dañoso.
\item Corren dos años de caducidad (usualmente el demandante se los toma).
\item \textbf{Primer momento de cuantificación: la demanda.} El art. 162 del CPACA exige como requisito de la demanda la estimación razonable de la cuantía: el demandante tiene que decir a cuánto asciende el valor de los perjuicios.
\item \textbf{Segundo momento de cuantificación: la sentencia.} El juez condena y tasa el perjuicio; las partes deben revisar que la tasación esté bien hecha.
\end{itemize}
Por eso la primera pregunta que hay que hacerse es: ¿para qué estoy calculando los perjuicios? ¿Para presentar la demanda, para verificar la condena o, como juez, para hacer la condena? La respuesta cambia la actualización, la tasación del perjuicio y las reglas que se aplican.

\begin{ejemplo}Gastos funerarios. Un soldado de 25 años muere; deja una cónyuge de 23 años y dos hijos de dos años y de un año. El día de la muerte la familia paga gastos funerarios por cinco millones de pesos. Esos cinco millones, pagados hace dos años, no valen cinco millones hoy, por lo que deben actualizarse para que las víctimas no pierdan dinero y sean realmente reparadas. Si se está en la demanda: IPC vigente al momento de la demanda ÷ IPC vigente al momento del pago, multiplicado por cinco millones. Si se está en la condena: IPC vigente al momento de la condena ÷ IPC vigente al momento del pago, multiplicado por cinco millones.\end{ejemplo}
\textbf{Cada gasto se actualiza por separado.} Si hubo varios gastos en fechas distintas (por ejemplo, los gastos funerarios pagados a los tres días y, a raíz de la muerte, consultas de la esposa con el psicólogo), cada uno debe actualizarse según la fecha en que se hizo. No se puede ser perezoso y actualizar todo desde una sola fecha: si no se actualiza cada gasto, el cliente pierde plata.

\subsection{Consejos prácticos para la actualización}
\begin{itemize}
\item \textbf{El resultado de la división siempre es positivo y mayor que uno.} Si da un número negativo o menor que uno, algo va mal: cinco millones de hace dos años no pueden convertirse en tres millones, porque el dinero no aumenta de valor; pierde valor.
\item \textbf{El número tiene que ser razonable.} Siempre será «uno y algo chiquitico» (1,1; 1,2; 1,5). Si la división da 7, algo va mal: jamás cinco millones se convierten en 35 millones en dos años. «Si le da un número muy grande, preocúpese, porque o está tumbando al cliente o se está tumbando usted.»
\item \textbf{El IPC se publica mensualmente.} El DANE tiene la obligación legal de publicar el IPC dentro de la primera semana de cada mes. Si hoy es 2 de octubre de 2026, todavía no está publicado el IPC de octubre (saldrá el 7 de octubre); por lo tanto, como índice final debe usarse el de septiembre, que es el último publicado.
\item \textbf{Aunque el bien cueste lo mismo, el dinero vale menos.} Si algo que costaba cinco millones hace dos años sigue costando cinco millones hoy, de todos modos hay que actualizar, porque los cinco millones de hace dos años no son los mismos de hoy: con un millón de pesos hace dos años se podía hacer un mercado mucho más completo que el que se puede hacer hoy con la misma suma. El servicio fúnebre puede costar un millón antes y un millón hoy, pero el millón de hace dos años no es el mismo millón de hoy.
\end{itemize}
\section{El lucro cesante}
\begin{cita}«Art. 1614. (…) Entiéndese (…) por lucro cesante, la ganancia o provecho que deja de reportarse a consecuencia de no haberse cumplido la obligación, o cumplido imperfectamente, o retardado su cumplimiento.» (Código Civil)\end{cita}
\subsection{Lucro cesante consolidado y futuro}
El momento hasta el cual se extiende el lucro cesante consolidado, y desde el cual arranca el futuro, depende del momento procesal en el que se hace la liquidación:

{\small\begin{xltabular}{\linewidth}{|X|X|X|}\hline
\rowcolor{navy!12}\textbf{Momento de la liquidación} & \textbf{Lucro cesante consolidado} & \textbf{Lucro cesante futuro} \\\hline\endhead
En la demanda & Desde los hechos hasta la presentación de la demanda & Desde la demanda hasta el vencimiento de la expectativa de vida \\\hline
En la sentencia & Desde los hechos hasta la sentencia & Desde la sentencia hasta el vencimiento de la expectativa de vida \\\hline
\end{xltabular}}
La distinción no es solo académica: el lucro cesante consolidado y el futuro tienen fórmulas distintas, porque los intereses que se aplican en uno y otro son diferentes.

\subsection{Fórmula del lucro cesante consolidado}
\formula{LCC = Ra × [ (1 + i)ⁿ − 1 ] ÷ i}
\begin{itemize}
\item \textbf{Ra:} renta mensual actualizada (el ingreso mensual actualizado de la víctima).
\item \textbf{i:} interés puro o técnico mensual = 0,004867. Parece una variable, pero no lo es: es una constante que siempre es igual, tanto en el lucro cesante consolidado como en el futuro.
\item \textbf{n:} número de meses que comprende el período a indemnizar. Como en la fórmula \emph{n} es una potencia, se calcula (1 + i) y se eleva a la \emph{n} (se necesita calculadora científica).
\end{itemize}
Por eso, realmente, la fórmula solo tiene dos variables: \emph{n} y \emph{Ra}.

\textbf{¿Cómo se determina \emph{n}?} En meses. Por regla general, si se esperó todo el término de caducidad para demandar, son dos años: el lucro cesante consolidado será de 24 meses y \emph{n} = 24. Si se liquida en la sentencia y han pasado, por ejemplo, 10 años, \emph{n} = 10 × 12 = 120. Si fueron 9 años y 5 meses, se redondea hasta obtener el exponente.

\textbf{¿Por qué la renta debe ser mensual?} Porque \emph{n} está en meses; si no se usa el ingreso mensual, la fórmula no cuadra.

\textbf{Orden de las operaciones:} (1) se calcula 1 + i; (2) se eleva a la potencia \emph{n}; (3) se le resta 1; (4) se divide por \emph{i}; (5) el resultado se multiplica por la renta mensual actualizada.

\subsection{Fórmula del lucro cesante futuro}
\formula{LCF = Ra × [ (1 + i)ⁿ − 1 ] ÷ [ i × (1 + i)ⁿ ]}
Tiene las mismas variables: \emph{Ra} sigue siendo la renta mensual actualizada, \emph{i} = 0,004867 y \emph{n} sigue siendo el número de meses del período a indemnizar. Lo que cambia es la fórmula y el período que comprende \emph{n}: desde la demanda (o la sentencia) hasta el vencimiento de la expectativa de vida.

\textbf{Orden de las operaciones:} (1) se calcula 1 + i y se eleva a la \emph{n}; (2) se le resta 1 (numerador); (3) se calcula de nuevo (1 + i)ⁿ y se multiplica por \emph{i} (denominador); (4) se divide el numerador por el denominador; (5) el resultado se multiplica por \emph{Ra}.

\begin{ejemplo}Supongamos una expectativa de vida de 30 años (dato inventado para el ejemplo). Si se liquida para la demanda: el consolidado es de 24 meses y el futuro de 28 × 12 = 336 meses. Si se liquida en una sentencia dictada a los 10 años: el consolidado es de 120 meses y el futuro de 20 × 12 = 240 meses.\end{ejemplo}
\subsection{Reglas jurisprudenciales comunes: la renta (Ra)}
1. \textbf{Incremento del 25 \% por prestaciones sociales.} Primero hay que preguntarse si la persona recibía prestaciones sociales, lo que ocurre cuando tiene contrato de trabajo o es empleado público (no así el contratista por prestación de servicios ni el independiente). Si las recibía, se suma un 25 \% a la renta.

\begin{ejemplo}El soldado ganaba tres millones de pesos mensuales y, como estaba nombrado mediante acto administrativo, recibía prestaciones sociales: 3.000.000 × 1,25 = 3.750.000. Si fuera contratista, no se suma el 25 \%.\end{ejemplo}
2. \textbf{Actualización de la renta.} Por la pérdida de valor del dinero, la renta debe actualizarse a la fecha de la presentación de la demanda o a la fecha de la sentencia, según el momento de la liquidación (IPC final ÷ IPC inicial × 3.750.000).

3. \textbf{Piso del salario mínimo.} Si después de actualizarla la renta es inferior al salario mínimo legal mensual vigente (SMLMV), se liquida con el SMLMV vigente a la fecha de la liquidación, porque se presume que nadie en Colombia devenga menos del salario mínimo.

4. \textbf{Presunción del salario mínimo cuando no se prueban los ingresos.} Si se acredita que la víctima ejercía una actividad económica productiva, pero no se prueba el valor de sus ingresos, se liquida con el salario mínimo (fue lo que ocurrió en el caso de alias «Bonais»).

\begin{ejemplo}El abogado privado injustamente de la libertad. En las profesiones liberales es frecuente que los demandantes se concentren en demostrar que la víctima era un abogado muy bueno, que estaba litigando, pero olviden probar el monto de sus ingresos. Resultado: se liquida con el salario mínimo. ¿Qué debió aportarse? Las declaraciones de renta del año anterior o de los últimos años, para sacar un promedio. Si la víctima era empleada, se aporta la certificación del último empleador sobre el salario o la remuneración que recibía (si la víctima directa murió, la familia debe aportarla).\end{ejemplo}
5. \textbf{Amas y amos de casa.} La jurisprudencia de unificación ha reconocido que ser ama o amo de casa equivale a un trabajo: la dedicación al hogar ocupa el mismo tiempo laboral. Se liquida con base en el salario mínimo.

\subsection{Reglas comunes: las víctimas indirectas}
Cuando la víctima directa no muere (por ejemplo, un soldado que gana tres millones y es privado injustamente de la libertad durante seis meses), el lucro cesante se le reconoce a ella, que era quien ejercía la actividad económica. ¿Qué pasa cuando quien ejercía la actividad productiva muere y reclaman sus familiares?

\begin{itemize}
\item \textbf{Solo pueden reclamar quienes tenían dependencia económica con la víctima directa}, y esa dependencia debe probarse en el proceso.
\item \textbf{Presunción a favor de los hijos menores de 25 años:} los hijos solo tienen que demostrar que son hijos y que eran menores de 25 años; se presume que dependían económicamente del padre.
\item \textbf{A la cónyuge no se le presume la dependencia económica:} debe demostrarla (pruebas de que el esposo ayudaba en la casa, de que recibía parte de lo que él ganaba, etc.). Si no la prueba, no puede reclamar lucro cesante.
\item \textbf{Cualquier persona puede probar dependencia económica}, sea o no familiar: la cónyuge, los padres, los sobrinos e incluso el compadre o el amigo con el que la víctima jugaba billar, si demuestra que mensualmente la víctima le pagaba algunos gastos o lo ayudaba con los pagos de la casa. Es un problema de prueba.
\end{itemize}
\textbf{La «bolsa» de dinero y los gastos personales.} El lucro cesante es una bolsa de dinero de la que se alimentaban varias personas. La víctima directa tenía gastos personales (se alimentaba, salía a comer) que, habiendo muerto o no, nunca habrían llegado a la familia. Por eso no se liquida con toda la bolsa:

\begin{itemize}
\item Si solo hay un demandante (por ejemplo, solo la cónyuge), los gastos personales son el 50 \%: se liquida sobre el 50 \% restante.
\item Si hay más demandantes (cónyuge e hijos), los gastos personales se reducen al 25 \%. Tiene lógica: entre más grande la familia, la persona gasta menos en sí misma y destina más dinero a la familia.
\item Después de descontar los gastos personales, la renta se divide: 50 \% para la cónyuge y 50 \% para los hijos (que se reparte entre ellos).
\end{itemize}
\begin{ejemplo}Distribución en el caso del soldado (renta de 3.750.000; cónyuge y dos hijos). Se descuenta el 25 \% de gastos personales; queda el 75 \%. La mitad, 37,5 \%, es para la cónyuge (1.406.250); el otro 37,5 \% se distribuye entre los dos hijos: 18,75 \% para cada uno (703.125).\end{ejemplo}
\fuente{En el ejemplo no se aplica la actualización de la renta, para simplificar las cifras.}
\subsection{Reglas comunes: el período indemnizable}
\begin{itemize}
\item \textbf{Consolidado:} desde la ocurrencia del daño hasta la presentación de la demanda o hasta la sentencia.
\item \textbf{Futuro:} desde la presentación de la demanda (o la sentencia) hasta la expectativa probable de vida.
\item \textbf{Acrecimiento:} liquidación por períodos.
\end{itemize}
\textbf{La expectativa de vida.} Se usa la tabla de la Superintendencia Financiera (Resolución 0110 del 22 de enero de 2014), con las tasas de mortalidad y expectativa de vida probable que se usan en materia de seguros y también en responsabilidad extracontractual. La tabla indica, por sexo y edad, cuántos años de vida probable le quedan a una persona.

\begin{ejemplo}En el caso del soldado (hombre de 25 años), la tabla arrojó una expectativa de 52,9 años; para la cónyuge (mujer de 23 años), 60,6 años.\end{ejemplo}
\textbf{Regla: siempre se toma la menor expectativa de vida.} El razonamiento es: si el soldado no hubiera pisado la mina, ¿cuándo habría muerto? Si habría muerto antes que su cónyuge, sus ingresos solo habrían llegado hasta los 52,9 años. Y si el resultado fuera al revés (el soldado con 60 y la cónyuge con 52,9), también se toma 52,9: si el soldado no hubiera muerto, de todos modos la cónyuge habría muerto antes y no se le puede dar algo que nunca habría devengado o aprovechado.

\textbf{Son presunciones que admiten prueba en contrario.} Tanto la regla de los 25 años de los hijos como las tablas de mortalidad son estimaciones:

\begin{itemize}
\item Un hijo con algún tipo de discapacidad, si la prueba, puede recibir lucro cesante no hasta los 25 años, sino según su expectativa de vida probable.
\item El demandado puede probar, por ejemplo, que el soldado tenía un cáncer muy agresivo y que el médico le había dicho que solo viviría cinco años; en ese caso se reduce el período \emph{n}.
\end{itemize}
\subsection{Ejercicio completo: liquidación con acrecimiento (para la demanda)}
\begin{ejemplo}Datos: soldado de 25 años (expectativa 52,9 años), cónyuge de 23 años (expectativa 60,6 años), hijo 1 de dos años e hijo 2 de un año. Renta: 3.750.000 (con prestaciones). Se liquida para la presentación de la demanda (24 meses después de los hechos).\end{ejemplo}
{\small\begin{xltabular}{\linewidth}{|X|X|X|X|X|}\hline
\rowcolor{navy!12}\textbf{Período} & \textbf{Meses (n)} & \textbf{Beneficiarios} & \textbf{Distribución de la renta} & \textbf{Ra por beneficiario} \\\hline\endhead
Lucro cesante consolidado & 24 & Cónyuge, hijo 1 e hijo 2 & 25 \% gastos personales; 37,5 \% cónyuge; 18,75 \% cada hijo & Cónyuge: 1.406.250 / cada hijo: 703.125 \\\hline
Lucro cesante futuro (hasta que el hijo 1 cumple 25 años) & 252 (= 276 − 24) & Cónyuge, hijo 1 e hijo 2 & Igual a la anterior & Cónyuge: 1.406.250 / cada hijo: 703.125 \\\hline
Lucro cesante futuro (hasta que el hijo 2 cumple 25 años) & 12 & Cónyuge e hijo 2 & 25 \% gastos personales; 37,5 \% cónyuge; 37,5 \% hijo 2 (acrecimiento) & Cónyuge: 1.406.250 / hijo 2: 1.406.250 \\\hline
Lucro cesante futuro (resto de la expectativa de vida) & 346,8 & Solo la cónyuge & 50 \% gastos personales; 50 \% cónyuge & Cónyuge: 1.875.000 \\\hline
\end{xltabular}}
Cálculos:

\begin{itemize}
\item Período total: 52,9 años × 12 = 634,8 meses.
\item El hijo 1 (dos años) recibe hasta los 25 años: 23 años × 12 = 276 meses. El hijo 2 (un año) recibe 12 meses más.
\item Último período: 634,8 − 24 − 252 − 12 = 346,8 meses. Verificación: 346,8 + 12 + 252 + 24 = 634,8.
\end{itemize}
\textbf{Operaciones:} en el consolidado se hacen tres liquidaciones (cónyuge, hijo 1 e hijo 2) con \emph{n} = 24. Luego se aplica la fórmula del lucro cesante futuro con \emph{n} = 252 para los tres. En los 12 meses siguientes la bolsa cambia: como hay un hijo menos, «la bolsa vuelve a crecer» y se le puede dar más al otro hijo (esto es el acrecimiento). Finalmente, cuando solo queda la cónyuge, los gastos personales vuelven a ser del 50 \% y a la cónyuge le corresponde el otro 50 \% (1.875.000), con \emph{n} = 346,8.

\begin{comentario}Como la bolsa de dinero se modifica a medida que avanza el tiempo, para verificar si el juez liquidó bien hay que hacer toda esta serie de cálculos.\end{comentario}
\section{Los perjuicios morales}
\subsection{Antecedentes: la irreparabilidad de los perjuicios morales}
Inicialmente, ni la jurisdicción ordinaria ni la contencioso administrativa reconocían perjuicios morales, por dos razones:

\begin{itemize}
\item \textbf{«Las lágrimas no se monetizan»} (o «las lágrimas no se monedean»). Era una teoría francesa: el dolor no se puede convertir en dinero.
\item \textbf{La imposibilidad de calcular su valor:} es imposible determinar cuánto sufrió una persona por la pérdida de un ser querido. Se trata de la imposibilidad de ordenar una reparación equivalente, debido a que «el dinero no serviría para reemplazar un bien que no tiene valor pecuniario».
\end{itemize}
\subsection{El fallo Villaveces (CSJ, 21 de julio de 1922)}
\begin{ejemplo}Los hechos. El señor Villaveces era viudo y visitaba periódicamente la tumba de su esposa. En una visita descubrió que habían exhumado el cadáver sin su consentimiento. Reclamó perjuicios morales: «¿Cómo así que ustedes retiran la tumba de mi esposa, la exhuman y no me piden ni permiso ni consentimiento?». Por primera vez, la Corte Suprema de Justicia reconoció los perjuicios morales.\end{ejemplo}
La reparación de los perjuicios morales se dedujo de los arts. 2341 y 2356 del C.C.:

\begin{cita}«Tanto se puede dañar a un individuo menoscabando su hacienda, como infligiéndole ofensa en su honra o en su dignidad personal o causándole dolor o molestia por obra de malicia o negligencia en el agente.»\end{cita}
\begin{cita}«Todo derecho lesionado requiere una reparación.»\end{cita}
\begin{cita}«Si en muchos casos es difícil determinar el quantum de la reparación, esa circunstancia no puede ser óbice para fijarlo aunque sea aproximadamente.»\end{cita}
\begin{definicion}\textbf{Perjuicios morales:} buscan reparar —en estricto sentido, compensar— la zozobra, el dolor y la angustia causados por el hecho dañoso.\end{definicion}
\textbf{Cómo superó la Corte las dos críticas:}

\begin{itemize}
\item \textbf{Compensar no es indemnizar.} Indemnizar es restituir las cosas al estado anterior al hecho dañoso: dejar a la persona en la misma situación en la que estaba, pagándole el valor en dinero de lo que perdió o dejó de ganar. Los perjuicios morales no indemnizan, porque no hay valor económico que reemplace la muerte de un ser querido; lo que hacen es compensar: mitigar, ayudar a una persona en dificultad a salir adelante. Nunca buscan que las cosas vuelvan al estado anterior, sino reducir la aflicción para que la persona pueda superar su estado de zozobra. El dolor, en efecto, no se monetiza, pero una ayuda económica puede ayudar a salir adelante.
\item \textbf{Fórmulas de equidad.} Es imposible saber a cuánto asciende el dolor, pero pueden usarse fórmulas de equidad y de arbitrio judicial para determinar una compensación económica suficiente.
\end{itemize}
\subsection{La analogía con el artículo 95 del Código Penal de 1936}
Inicialmente, tanto la Corte Suprema como el Consejo de Estado aplicaron, como tope para tasar los perjuicios morales, el artículo 95 del Código Penal de 1936:

\begin{cita}«Cuando no fuere fácil o posible avaluar económicamente el daño moral, ocasionado por el delito, el Juez puede fijar prudencialmente la indemnización que corresponda al ofendido, hasta la suma de dos mil pesos.»\end{cita}
El tope era, entonces, de 2.000 pesos.

\subsection{El abandono de la analogía en la Corte Suprema (CSJ, 27 de septiembre de 1974)}
\begin{cita}«Esta limitación sólo tiene cabida en los precisos eventos de regulación del "daño moral ocasionado por el delito"; además, el precepto está dirigido a los jueces penales y no a los de otras jurisdicciones.»\end{cita}
Razones: el código ya había sido derogado (¿cómo seguir aplicando un código derogado?) y era una norma penal que no tenía sentido aplicar a la responsabilidad civil extracontractual.

\textbf{Los topes jurisprudenciales.} Desde 1974 y hasta hace poco, la Corte Suprema fijó, mediante sentencias, sumas sin un criterio objetivo de actualización: hacia 2007 el tope era de alrededor de 14 millones de pesos; luego se reconocieron 44 millones, 75 millones hacia 2018, y se llegó hasta unos 78 millones. Cada cierto tiempo se expedía una sentencia actualizando el tope; era un máximo (podía reconocerse menos) y nunca se actualizaba automáticamente. Excepcionalmente, se llegaron a reconocer 150 millones.

\subsection{Los gramos oro y el salario mínimo en el Consejo de Estado}
\begin{itemize}
\item \textbf{CE, 9 de febrero de 1978:} abandono de la analogía del art. 95 del Código Penal de 1936, por la desvalorización del monto señalado en la norma y por su aplicación limitada al derecho penal. Se adopta la tasación en gramos oro, con un tope de 1.000 gramos oro (equivalente a los 2.000 pesos). La ventaja era que el gramo oro sube de valor con el tiempo, de modo que la indemnización no pierde valor.
\item \textbf{CE, Sección Tercera, 6 de septiembre de 2001 (ponencia del doctor Alier Hernández):} se abandonan los gramos oro —una figura «anacrónica», propia «del viejo oeste»— y se pasa al salario mínimo, que se actualiza automáticamente y es más fácil, claro y moderno. El tope de 1.000 gramos oro se reemplazó por un tope de \textbf{100 SMLMV}, que es el que principalmente se aplica hoy.
\end{itemize}
\begin{comentario}La Corte Suprema adopta el salario mínimo. En una sentencia reciente, con ponencia del magistrado Octavio Tejeiro, la Corte Suprema de Justicia se acercó a la posición del Consejo de Estado y estableció como tope de los perjuicios morales 100 salarios mínimos. Es un cambio muy importante porque: (i) no tenía justificación usar sumas fijas que no se actualizaban; y (ii) rompe una lógica «nefasta»: el mismo daño generaba una indemnización distinta según lo causara un particular o una entidad estatal (por ejemplo, si un paciente con dolor abdominal moría en una clínica privada, cada familiar recibía unos 74 millones; si moría en un hospital público, recibía 100 salarios mínimos, mucho más). La equivalencia es más justa y equitativa.\end{comentario}
\subsection{Finalidad de los perjuicios morales}
\begin{itemize}
\item \textbf{Satisfactoria o compensatoria, no reparatoria} (CSJ, 27 de septiembre de 1974 / CE, S3, 20 de abril de 2005): no dejan indemne a la víctima ni devuelven las cosas al estado anterior.
\item \textbf{No buscan el enriquecimiento de la víctima} (CSJ, 21 de febrero de 2018).
\item \textbf{No tienen finalidad sancionatoria ni son daños punitivos} (Corte IDH, caso Velásquez Rodríguez): no buscan castigar al demandado.
\end{itemize}
\begin{comentario}«No hay muerto malo.» En la responsabilidad extracontractual, cuando muere alguien, no solo aparecen la esposa y los hijos: aparecen tíos, tatarabuelos y primos terceros, todos describiendo a la víctima como el familiar ejemplar; muchas demandas «parecen más herencias». Consejo práctico: una buena puerta de entrada con el juez es ser moderado con los perjuicios. Una demanda con la esposa, los hijos, los hermanos, los padres e incluso los abuelos es completamente razonable; cuando aparecen primos segundos y terceros, el juez intuye que no se busca una reparación. Además, si se pierde el proceso, la condena en costas es proporcional a la cuantía de los perjuicios reclamados.\end{comentario}
\subsection{Las sentencias de unificación de 2014 y 2021}
El 28 de agosto de 2014 el Consejo de Estado expidió nueve sentencias de unificación sobre perjuicios inmateriales en tres hechos dañosos —muerte, lesiones personales y privación injusta de la libertad— y unificó dos asuntos: la cuantificación y la prueba de los perjuicios morales. En privación injusta de la libertad, las reglas de 2014 fueron modificadas en 2021 mediante una nueva unificación, porque tenían varios errores que hacían injustas las tasaciones.

Sentencias de unificación de perjuicios morales:

\begin{itemize}
\item Muerte: CE, S3, 26251, 28 de agosto de 2014 / CE, S3, 27709, 28 de agosto de 2014.
\item Lesiones: CE, S3, 31172, 28 de agosto de 2014.
\item Privación injusta de la libertad: CE, S3, 46681, 29 de noviembre de 2021.
\end{itemize}
\subsection{Perjuicios morales en casos de muerte (26251 y 27709)}
Dependen de la cercanía del demandante con la víctima directa: entre más cercano, mayor el perjuicio; entre más lejano, menor.

{\small\begin{xltabular}{\linewidth}{|X|X|X|X|}\hline
\rowcolor{navy!12}\textbf{Nivel} & \textbf{Relación con la víctima directa} & \textbf{Porcentaje} & \textbf{SMLMV} \\\hline\endhead
Nivel 1 & Relaciones afectivas conyugales y paterno-filiales (padres, hijos, cónyuge, compañeros permanentes) & 100 \% & 100 \\\hline
Nivel 2 & Relación afectiva del 2.º de consanguinidad o civil (abuelos, hermanos y nietos) & 50 \% & 50 \\\hline
Nivel 3 & Relación afectiva del 3.º de consanguinidad o civil & 35 \% & 35 \\\hline
Nivel 4 & Relación afectiva del 4.º de consanguinidad o civil & 25 \% & 25 \\\hline
Nivel 5 & Relaciones afectivas no familiares (terceros damnificados) & 15 \% & 15 \\\hline
\end{xltabular}}
\textbf{¿Quién puede reclamar?} Cualquier persona que acredite haber sufrido por la muerte de la víctima directa; lo que cambia es la cuantía, según la cercanía.

\textbf{Prueba:} en los niveles 1 y 2 los perjuicios morales se presumen a partir de la prueba del parentesco: con los registros civiles se tiene por probado el dolor y no hay que probar nada más. En los demás niveles no opera la presunción: hay que demostrarle al juez la cercanía con la víctima y el dolor que su muerte produjo.

\textbf{Nota 2 (26251):}

\begin{cita}«En casos excepcionales, como los de graves violaciones a los derechos humanos, entre otros, podrá otorgarse una indemnización mayor de la señalada en todos los eventos anteriores, cuando existan circunstancias debidamente probadas de una mayor intensidad y gravedad del daño moral, sin que en tales casos el monto total de la indemnización pueda superar el triple de los montos indemnizatorios antes señalados. Este quantum deberá motivarse por el juez y ser proporcional a la intensidad del daño.»\end{cita}
\begin{comentario}¿Función sancionatoria encubierta? El tope es de 100 salarios, pero cuando el hecho dañoso proviene de una grave violación de derechos humanos puede triplicarse hasta 300. Esto tiene un componente sancionatorio: la muerte de una persona, independientemente de cómo se cause, debería producir el mismo dolor; si lo que permite incrementar el perjuicio no es la intensidad del dolor sino la gravedad de la conducta del demandado, se está sancionando. No es contrario a la igualdad (se aplica a todos los que estén en esa situación) ni supone excusar las violaciones de derechos humanos, pero sí tiene «un tinte sancionador», en contravía de lo dicho por la Corte Interamericana.\end{comentario}
\subsection{Perjuicios morales en casos de lesiones (31172)}
Hay dos variables: (i) el grado de cercanía con la víctima directa (los mismos niveles, con la diferencia de que en el nivel 1 se incluye también la víctima directa, que en las lesiones sigue viva y puede reclamar); y (ii) la gravedad de la lesión, determinada por el porcentaje de pérdida de capacidad laboral.

{\small\begin{xltabular}{\linewidth}{|X|X|X|X|X|X|}\hline
\rowcolor{navy!12}\textbf{Gravedad de la lesión} & \textbf{Nivel 1 (víctima directa y relaciones conyugales y paterno-filiales)} & \textbf{Nivel 2} & \textbf{Nivel 3} & \textbf{Nivel 4} & \textbf{Nivel 5} \\\hline\endhead
Igual o superior al 50 \% & 100 & 50 & 35 & 25 & 15 \\\hline
Igual o superior al 40 \% e inferior al 50 \% & 80 & 40 & 28 & 20 & 12 \\\hline
Igual o superior al 30 \% e inferior al 40 \% & 60 & 30 & 21 & 15 & 9 \\\hline
Igual o superior al 20 \% e inferior al 30 \% & 40 & 20 & 14 & 10 & 6 \\\hline
Igual o superior al 10 \% e inferior al 20 \% & 20 & 10 & 7 & 5 & 3 \\\hline
Igual o superior al 1 \% e inferior al 10 \% & 10 & 5 & 3,5 & 2,5 & 1,5 \\\hline
\end{xltabular}}
\fuente{Valores en SMLMV (Gráfico No. 2 de la sentencia).}
Para saber el valor, se ubica qué tan lejano es el demandante de la víctima y cuál fue el porcentaje de pérdida de capacidad, y se busca en la tabla. \textbf{Prueba:} la sentencia no fija regla expresa sobre la presunción derivada del parentesco, pero en el caso concreto la aplica para los niveles 1 y 2; opera la misma lógica que en muerte: en los niveles 1 y 2 se presume el dolor con la prueba del parentesco; en los demás, debe probarse.

\textbf{¿Cómo se prueba el dolor fuera de los niveles 1 y 2?} Con el testimonio del psicólogo que atendió a la persona (por ejemplo, para acreditar que el amigo de la víctima asistió a terapia a raíz de su muerte); incluso la prueba testimonial podría ser suficiente. El problema es la cuantía: un amigo, tanto en muerte como en lesiones, está en el nivel de tercero damnificado (15 salarios). La prueba psicológica puede ayudar a justificar un valor más alto cuando había una relación muy cercana.

\textbf{Relaciones de crianza.} En el nivel 1 también se encuentran los hijos y padres de crianza: en responsabilidad extracontractual, las relaciones de crianza son equivalentes a las biológicas.

\textbf{Autodeclaraciones.} ¿Puede un amigo acreditar la amistad con una declaración juramentada? No: las autodeclaraciones no sirven: ir a una notaría y rendir una declaración extrajuicio diciendo «yo era el mejor amigo» es certificar uno mismo lo que quiere probar, y la contraparte debe poder contradecir la prueba. De lo contrario, cualquiera declararía extrajuicio que ganaba 800 millones de pesos mensuales.

\textbf{El cargo de la víctima no altera los perjuicios morales.} Los morales son los mismos si muere el Presidente de la República o un desempleado: el dolor no depende del cargo de la persona, sino de la cercanía con ella. En cambio, en el lucro cesante sí son relevantes los ingresos (la variable \emph{Ra} depende del cargo).

\subsection{Perjuicios morales en casos de privación injusta de la libertad (46681, 2021)}
\begin{itemize}
\item Se presumen los perjuicios morales para la víctima directa.
\item A partir de la prueba de la relación o parentesco con la víctima directa, se presumen los perjuicios morales para los cónyuges, compañeros/as permanentes y parientes hasta el primer grado de consanguinidad.
\item \textbf{Cuantificación para la víctima directa:}
\begin{itemize}
\item Un mes o menos: 5 SMLMV.
\item Si la privación dura más de un mes: PM = (número de meses × 5 SMLMV) + (fracción adicional de días × 0,166 SMLMV).
\item Tope: 100 SMLMV; excepcionalmente, 300 SMLMV.
\item Descuento en casos de detención domiciliaria: hasta el 50 \%.
\end{itemize}
\item \textbf{Víctimas indirectas cercanas} (padres, hijos, cónyuges y compañeros permanentes): hasta el 50 \% de lo que corresponda a la víctima directa.
\item \textbf{Otras víctimas indirectas:} hasta el 35 \% de lo que corresponda a la víctima directa.
\end{itemize}
\begin{comentario}Para la tasación en privación injusta deben consultarse la sentencia de unificación de 2021 (perjuicios morales) y la de 2019 (perjuicios materiales). En estas reglas ya no es relevante, para los morales, el ingreso que tenía la persona; eso era antes y cambió.\end{comentario}
\section{El daño a la salud}
\subsection{Antecedentes: del perjuicio fisiológico al daño a la vida de relación}
Antes de hablarse de daño a la salud, solo existían los tres perjuicios clásicos: daño emergente, lucro cesante y perjuicios morales.

\begin{ejemplo}Freddy deja de jugar fútbol. Freddy no puede reclamar lucro cesante por el fútbol porque no es futbolista profesional. Pero si a raíz de un accidente de tránsito deja de jugar fútbol, que era su pasión y su pasatiempo de fin de semana, ¿qué perjuicio sufre? La jurisprudencia dijo que no es en estricto sentido un perjuicio moral (no es dolor ni zozobra), sino la afectación del ejercicio de actividades placenteras de la vida.\end{ejemplo}
\begin{itemize}
\item \textbf{Perjuicio fisiológico:} derivado de una lesión física o corporal; afecta los placeres de la vida.
\item \textbf{Daño a la vida de relación:} la jurisprudencia consideró que el concepto de perjuicio fisiológico era demasiado estrecho porque lo ataba a lesiones físicas o corporales, y no siempre hay una lesión corporal que afecte los placeres de la vida o la relación con el mundo exterior: puede tratarse de afectaciones emocionales que afectan la autoestima. Va más allá de las lesiones físicas o corporales; comprende las actividades cotidianas y las relaciones con otros seres y los actos individuales.
\end{itemize}
\begin{ejemplo}La cicatriz. Una persona es atropellada al salir de clase y, tras la operación, queda con una cicatriz. No pierde capacidad laboral, pero la cicatriz afecta su autoestima: le encantaba ir a la playa y asolearse, y ya no puede hacerlo porque siente vergüenza de su cuerpo. Además de los perjuicios morales por el dolor del accidente, puede pedir daño a la vida de relación por la afectación de su relación con el mundo exterior.\end{ejemplo}
\subsection{El daño a la salud (unificación de 2014)}
En 2014 se abandona el daño a la vida de relación y se adopta el concepto actual de daño a la salud (CE, S3, 31170, 28832 y 28804, todas del 28 de agosto de 2014).

\begin{definicion}\textbf{Daño a la salud:} según las sentencias de unificación, es el perjuicio inmaterial que se ocasiona por lesiones psicofísicas que sufre la víctima directa y que no se encuentran inmersas dentro del perjuicio moral. Solo lo sufre la víctima directa.\end{definicion}
{\small\begin{xltabular}{\linewidth}{|X|X|}\hline
\rowcolor{navy!12}\textbf{Gravedad de la lesión} & \textbf{Víctima directa} \\\hline\endhead
Igual o superior al 50 \% & 100 SMMLV \\\hline
Igual o superior al 40 \% e inferior al 50 \% & 80 SMMLV \\\hline
Igual o superior al 30 \% e inferior al 40 \% & 60 SMMLV \\\hline
Igual o superior al 20 \% e inferior al 30 \% & 40 SMMLV \\\hline
Igual o superior al 10 \% e inferior al 20 \% & 20 SMMLV \\\hline
Igual o superior al 1 \% e inferior al 10 \% & 10 SMMLV \\\hline
\end{xltabular}}
\begin{itemize}
\item En casos de extrema gravedad y excepcionales se podrá aumentar hasta 400 SMMLV, siempre que esté debidamente motivado.
\item En casos de daño a la salud por lesiones temporales, para su tasación debe establecerse un parangón con el monto máximo que se otorgaría en caso de lesiones similares a aquellas objeto de reparación, pero de carácter permanente y, a partir de allí, determinar la indemnización en función del período durante el cual, de conformidad con el acervo probatorio, se manifestaron las lesiones a indemnizar (28832).
\item No solo procede en casos de incapacidad: hay otras variables, de carácter cualitativo (28804). Si no hubo pérdida de capacidad laboral, el juez debe acudir a factores cualitativos que le permitan tasar el daño a la salud con arbitrio judicial, pero siempre debe tratarse de una lesión psicofísica que produzca algo distinto del daño moral.
\end{itemize}
\begin{comentario}Crítica: el daño a la salud «es un enredo», porque su definición no es clara ni precisa. Primero, si se trata de una lesión psíquica, ¿qué lesiones psíquicas no quedan dentro de los perjuicios morales? Es un problema conceptual imposible de escindir. Segundo, si es una lesión física, para tasarlo se usa el grado de incapacidad laboral, lo que lo confunde con el lucro cesante. La sensación es que el daño a la salud es redundante con los perjuicios morales o con el lucro cesante: cuando hay pérdida de capacidad laboral, se indemniza dos veces lo mismo. Al parecer, los jueces contencioso administrativos consideraron insuficientes los topes y se inventaron estas tipologías para dar más dinero («démosle además de 100 salarios por morales, otros 100 por daño a la salud»). Sería más fácil volver a un esquema con solo daño emergente, lucro cesante y perjuicios morales, y decir abiertamente que el perjuicio moral puede ir hasta 300, 400 o 500 salarios cuando se pruebe algo extremo.\end{comentario}
\begin{ejemplo}Privación injusta y estrés postraumático. A una persona privada injustamente de la libertad durante 20 meses o más se le reconocen 100 salarios por perjuicios morales (el tope). Si además aporta un dictamen psiquiátrico que muestra estrés postraumático, se le reconocen otros 100 (o 50, o lo que estime el juez) por daño a la salud. ¿Hasta qué punto ese estrés postraumático o las citas al psiquiatra no son simplemente la prueba de un perjuicio moral muy intenso, pero no de un perjuicio distinto?\end{ejemplo}
\begin{comentario}Precisiones. ¿Ya no se habla de daño a la vida de relación en lo contencioso administrativo? Respuesta: se subsume en los perjuicios morales y en el daño a la salud (principalmente en los morales). Por otra parte, la indemnización objetiva que pagan las ARL por pérdida de capacidad laboral se parece al daño a la salud: hay un puente común entre ambos escenarios.\end{comentario}
\section{Daños por la vulneración de bienes constitucional y convencionalmente protegidos}
\fuente{Sentencia de unificación: CE, S3, 32988, 28 de agosto de 2014.}
\subsection{Origen: el Sistema Interamericano}
A principios de los años 2000 muchas víctimas del conflicto armado acudieron al Sistema Interamericano: habían ganado su proceso ante la jurisdicción contenciosa (por ejemplo, por un falso positivo o una desaparición causada por el Estado), pero consideraban que esa indemnización no era suficiente porque no garantizaba sus derechos a la verdad, la justicia y la no repetición. La Corte Interamericana condenó al Estado colombiano, principalmente con condenas de carácter no pecuniario (por ejemplo, construir un monumento en honor a las víctimas para que los hechos no se repitan, o publicar la sentencia y hacer un libro a partir de ella). Ante esto, el Consejo de Estado creó una tercera tipología de perjuicio inmaterial diseñada para esos casos, con el fin de evitar que las víctimas demandaran luego al Estado ante el Sistema Interamericano.

\subsection{Reglas}
\begin{itemize}
\item \textbf{A solicitud de parte o de oficio:} es la única clase de perjuicio que el juez puede reconocer de oficio, aunque no se haya pedido en la demanda.
\item \textbf{Se repara principalmente a través de medidas de carácter no pecuniario:} actos de conmemoración, disculpas públicas, monumentos, difusión y publicación de la sentencia, etc.
\item \textbf{Indemnización excepcional:} en casos excepcionales cuya reparación integral, a consideración del juez, no sea suficiente, pertinente, oportuna o posible, podrá otorgarse una indemnización, única y exclusivamente a la víctima directa, mediante una medida pecuniaria de hasta 100 SMLMV, si fuere el caso, siempre y cuando la indemnización no hubiere sido reconocida con fundamento en el daño a la salud.
\item En caso de ordenarse esa indemnización excepcional, no debe estar comprendida dentro de los perjuicios materiales e inmateriales ya reconocidos.
\end{itemize}
\begin{comentario}Crítica: el estándar para la medida pecuniaria es muy difícil de argumentar. En un falso positivo en el que ya se triplicaron los morales hasta 300 salarios, se reconoció daño a la salud y lucro cesante, y se ordenó publicar la sentencia en la página del Ejército, que todos los soldados la leyeran y que se hiciera un acto de conmemoración y disculpas, ¿cómo justificar además 100 salarios? Es una carga argumentativa que el juez debe cumplir, pero hay jueces de reparación directa «súper flexibles» que dicen que todo es insuficiente sin justificarlo. La tendencia es que la excepción se está volviendo la regla general: ante cada grave violación de derechos humanos se triplican los morales, se dan medidas no pecuniarias, medidas pecuniarias y, además, daño a la salud. Es una manifestación del «discurso del constitucionalismo» en la responsabilidad.\end{comentario}
